% =====================================================================
%  BAB II TINJAUAN PUSTAKA
% =====================================================================
\chapter{TINJAUAN PUSTAKA}
\label{bab:tinjauan}

\petunjuk{Analisis kritis terhadap literatur, bukan rangkuman sumber (Pedoman
TA 7.2). \tahap{\WAJIB}{\WAJIB{} (final)}}

\section{Landasan Teori}
\petunjuk{Konsep fundamental yang dipakai: definisi, prinsip kerja algoritma,
arsitektur, paradigma. Proporsional dengan kebutuhan penelitian; tambahkan
subbab 2.1.1, 2.1.2, dst.\ sesuai kebutuhan.}

\isian{Tulis uraian dalam bentuk paragraf.}

% \subsection{Konsep A}
% \subsection{Konsep B}

\section{Penelitian Terdahulu (\textit{State of the Art})}
\petunjuk{Bahas penelitian sebelumnya secara naratif dan kritis, lalu ringkas
dalam tabel perbandingan seperti contoh Tabel~\ref{tab:penelitian-terdahulu}
(pedoman 7.2.2). Tabel hanya ringkasan; analisis tetap ditulis dalam
paragraf.}

\isian{Tulis uraian naratif dan kritis, lalu rujuk
Tabel~\ref{tab:penelitian-terdahulu}.}

\begin{table}[H]
\caption{Perbandingan Penelitian Terdahulu}
\label{tab:penelitian-terdahulu}
\renewcommand{\arraystretch}{1.5}
\begin{tabularx}{\linewidth}{|>{\raggedright\arraybackslash}p{2.1cm}|X|X|X|}
\hline
\rowcolor{bgkepala}
\textbf{Peneliti (tahun)} & \textbf{Metode dan dataset} &
\textbf{Hasil utama} & \textbf{Keterbatasan} \\ \hline
 & & & \\ \hline
 & & & \\ \hline
 & & & \\ \hline
 & & & \\ \hline
\end{tabularx}
\end{table}

\isian{Paragraf ringkasan kesenjangan: apa yang belum dijawab penelitian
terdahulu dan bagaimana penelitian ini mengisinya.}

\section{Kerangka Pemikiran}
\petunjuk{Alur logika yang menjembatani teori dan metodologi; sisipkan
diagram blok atau \textit{flowchart}.}

\isian{Uraikan kerangka pemikiran dan rujuk
Gambar~\ref{fig:kerangka}.}

\begin{figure}[H]
\centering
% Ganti kotak berikut dengan: \includegraphics[width=0.9\linewidth]{gambar/kerangka-pemikiran.png}
\fbox{\parbox[c][4cm][c]{0.85\linewidth}{\centering
  \isian{Sisipkan diagram kerangka pemikiran}}}
\caption{Kerangka Pemikiran Penelitian}
\label{fig:kerangka}
\end{figure}
